% Generated draft: requires mathematical and source review.
\documentclass[11pt]{article}
\usepackage[margin=25mm]{geometry}
\usepackage{fontspec}
\setmainfont{Latin Modern Roman}
\usepackage{amsmath,amssymb,mathtools,unicode-math}
\setmathfont{Latin Modern Math}
\usepackage{enumitem,needspace,xurl,hyperref}
\setlength{\parindent}{0pt}
\setlength{\parskip}{5pt}
\setlength{\emergencystretch}{3em}
\newcommand{\N}{\mathbb{N}}
\newcommand{\Z}{\mathbb{Z}}
\newcommand{\Q}{\mathbb{Q}}
\newcommand{\R}{\mathbb{R}}
\newcommand{\C}{\mathbb{C}}
\newcommand{\F}{\mathbb{F}}
\newcommand{\A}{\mathbb{A}}
\newcommand{\PP}{\mathbb{P}}
\newcommand{\E}{\mathbb{E}}
\newcommand{\eps}{\varepsilon}
\newcommand{\dd}{\,\mathrm{d}}
\begin{document}
\section*{UnsolvedMath ID: 20000238}
\textbf{Source problem code:} catalogue entry
\subsection{Definitions and mathematical statement}
Hailong Dao's main question concerns \(U=\operatorname{Spec}(R)\setminus\{\mathfrak m\}\), where \((R,\mathfrak m)\) is a local ring of dimension \(d\): what is the obstruction theory for splitting vector bundles on \(U\), especially those of rank \(d-1\)? The source does not specify regularity, residue field, or whether splitting means an \(\mathcal O_U\)-summand or a fuller decomposition. Satya Mandal separately asks for obstructions to a sheaf-surjection \(E\twoheadrightarrow\mathcal O_X\) on a general scheme. On a non-affine scheme such a surjection need not admit a global section, so existence of a covector and splitting of a given covector are distinct problems. Aravind Asok asks for obstruction theory for general module classes beyond projective modules; Dao asks for an Euler class group for reflexive modules over a regular local ring; another question asks for the role of \(\A^1\)-homotopy theory outside the affine case.

The attached report treats the following restricted criterion. Let \(R\) be a Noetherian normal local domain of dimension at least two, \(M\) a finite reflexive \(R\)-module locally free on \(U\), and \(E=\widetilde M|_U\). It asserts that a covector \(q:E\to\mathcal O_U\) corresponds to \(\varphi:M\to R\), and that
\[
q\text{ is sheaf-surjective}\quad\Longleftrightarrow\quad
\operatorname{length}_R(R/\varphi(M))<\infty.
\]
For such a surjection put \(K=\ker(q)\) and \(N=\ker(\varphi)\). The reported connecting construction gives a canonical injection
\[
R/\varphi(M)\hookrightarrow H^1(U,K)\cong H^2_{\mathfrak m}(N).
\]
The connecting image \(\delta(1)\) of \(1\in H^0(U,\mathcal O_U)=R\) vanishes exactly when the given \(q\) has a global right inverse. This is a criterion for splitting a specified surjection, not a complete Euler obstruction theory for all bundles.

Its explicit candidate counterexample takes
\[
R=k[x,y,z]_{(x,y,z)},\qquad
K=\ker\bigl(\mathcal O_U^3\xrightarrow{(x,y,z)}\mathcal O_U\bigr).
\]
The report asserts that \(K\) has rank two, \([K]=2[\mathcal O_U]\) in \(K_0(U)\), and every positive-degree ordinary algebraic Chern class in \(\operatorname{CH}^*(U)\) vanishes, while
\[
H^1(U,K)\cong k,\qquad K\text{ has no trivial line summand}.
\]
Thus the reported obstruction is invisible to those specified \(K_0\) and ordinary Chow-theoretic data; it does not assert vanishing of every possible refined Euler invariant.

The original remarks must also be retained. Paul Balmer suggests the Jouanolou affine-replacement trick. Asok writes \(\nu_r(X)\) for rank-\(r\) vector-bundle isomorphism classes and describes a non-affine \(Y\) with an \(\A^1\)-weak equivalence \(\A^n\to Y\) but large \(\nu_r(Y)\), warning that this set is not generally an \(\A^1\)-homotopy invariant. For a finite-dimensional \(k\)-vector space \(V\), let \(H\subset\PP(V)\times\PP(V^*)\) be the incidence hypersurface. The source calls
\[
p_1:(\PP(V)\times\PP(V^*))\setminus H\longrightarrow\PP(V)
\]
an \(\A^1\)-weak equivalence and asks whether \(p_1^*\) is surjective on vector-bundle isomorphism classes. Kirsten Wickelgren asks about Nisnevich/\'etale topology and realization into equivariant topological vector bundles on manifolds. These are related questions rather than consequences of the local criterion.

The report adds that an affine replacement does not remove descent: for an fpqc cover \(p:T\to X\), a section of a pulled-back bundle descends precisely when its two pullbacks to \(T\times_XT\) agree, as expressed by the usual equalizer of section sets. In particular a splitting after pullback must satisfy this compatibility to give a splitting on \(X\). The summary does not give a general representability or Euler-class construction resolving the original questions.
\subsection{Short English statement}
Ask for splitting obstructions on punctured spectra and more general non-affine schemes. The report proposes a local-cohomology obstruction for a given surjection and a rank-two example missed by ordinary Chern classes, while retaining the source's separate homotopy and descent questions.
\subsection{Sources}
\begin{itemize}
\item AIM workshop ``Projective modules and \(A^1\)-homotopy theory,'' May 2014, Question 1: \url{https://aimath.org/pastworkshops/projectiveA1problems.pdf}. The extracted remarks carry the date May 24, 2014 and the name Sarang Sane; no PDF was accessed here.
\item Local canonical registry \texttt{aim-algebraic-geometry-notes.json}, \texttt{notes[237]}, tag \texttt{question}; permanent UnsolvedMath ID \(20000238\), source code \texttt{catalogue entry}. Main and related questions are attributed to Hailong Dao, Satya Mandal, and Aravind Asok; remarks to Paul Balmer, Aravind Asok, and Kirsten Wickelgren. The source has substantial OCR corruption, including control-character artifacts, and no clean statement.
\item Supplied research attempt 1 proposing the finite-colength/local-cohomology criterion and explicit punctured-threefold example. No full authoritative bibliographic references are supplied; the original literature field is empty. The report is machine generated, not peer reviewed, and requires independent verification before citation.
\end{itemize}
\end{document}
